\chapter{系统维护与数据安全}

\section{系统日常维护与健康监测}

为了保证 LingoBridge 中文学习平台在长期运行过程中的高性能与高可用性，系统管理员与运维团队必须实施规范化的日常维护与健康监测计划。系统基于微服务与容器化理念进行架构设计，其底层模块的协作与维护遵循了计算机软件设计描述规范标准\cite{ieee1016}。在日常监测中，运维人员应当重点监控后端 Node.js 应用进程的 CPU 与内存占用率，防止由于高并发或内存泄漏引起的进程宕机。系统在服务端专门预留了健康检查接口（/api/v1/health），通过该接口可以实时获取当前数据库的连接状态、缓存命中率以及服务器物理磁盘的剩余空间。

除了被动的监控告警，运维团队应当每周进行一次系统日志的审计与分析。系统日志详细记录了用户的访问历史、异常报错以及长连接信令的延迟波动。通过对日志数据的深度挖掘，可以及时发现并排除潜在的弱网通信瓶颈或接口响应延迟。此外，随着多班级试点教学的持续进行，学生日常提交的音频录音数据量会呈线性增长。为了避免物理磁盘被撑爆，系统需配置定期的无用文件清理任务。在获取教学管理方授权的前提下，运维人员可使用自动化脚本，对超过一学期且已被批改完成的录音文件进行归档冷存储处理，或者对系统测试产生的模拟垃圾数据进行硬删除，从而保证服务器物理磁盘的健康空间比率。

\section{数据备份与容灾策略}

在数字化教育系统中，课程课件、学生作业以及录音数据是极其珍贵的教学资产，其完整性直接关系到学校试点的成败与教学研究的连续性。因此，系统确立了严格的数据备份与容灾保障体系。在物理服务器层面，数据库层应当开启定期自动备份任务。系统每天凌晨会自动执行数据库的导出操作，将当前的所有账号数据、课程关系以及作业提交轨迹打包为加密的 SQL 压缩包，并同步上传至异地备份云服务器中，实现数据的异地冗余保护。

在文件存储层面，对于教师上传的 PDF 课件与学生提交的跟读音频，系统在试点阶段采用本地磁盘存储与对象存储相结合的双轨备份策略。所有上传的媒体资源会在本地写入的同时，异步触发一条镜像上传任务，将文件备份至符合跨境合规的对象存储（OSS）中。一旦发生本地物理磁盘损坏或节点崩溃，系统可以秒级切换文件读取路由，直接从云端对象存储中拉取课件与音频资源，确保教学行为在灾难发生时依然不发生中断。运维人员还需每月举行一次数据恢复性演练，通过在沙盒环境中还原备份的 SQL 数据与文件对象，检验备份数据的可用性与恢复响应速度。

\section{数据安全与隐私合规}

作为一款面向国际高校留学生的中文教学辅助平台，数据的安全性与隐私保护是系统设计的最高准则。系统在开发全生命周期中均落实了最小化采集与精细化访问控制原则，切实保障用户个人数据的隐私权益。在访问控制方面，系统在后端接口层部署了基于角色的访问控制机制。教师仅能查询和批改自己名下关联班级的学生录音，无权跨班级或跨学校调取其他数据；学生仅能查看与自己账号绑定的个人学习进度和作业反馈；管理员虽然拥有系统级权限，但其所有管理操作均会触发系统审计日志，实现操作留痕与防抵赖。

在隐私合规方面，平台坚决不采集任何与教学行为无关的敏感生物特征或地理位置数据。用户密码在前端传输至后端后，会立即经过盐值混淆与哈希算法进行单向加密存储，即使数据库发生泄露，也无法反推出用户的明文密码。此外，为了符合全球主流的隐私合规标准，系统赋予了学生绝对的个人数据删除权。学生可以随时在个人控制台中对其历史录音进行硬删除。一旦学生触发删除指令，系统会立即从本地物理磁盘和云端对象存储中彻底抹除该音频文件，并清除数据库中与之相关联的所有行为轨迹，从根本上杜绝了学生隐私泄露的隐患。为了系统性梳理各项运维操作的执行周期与责任边界，表\ref{tab:maintenance_schedule}对日常的维护动作进行了明确界定。

\begin{table}[htbp]
  \centering
  \caption{LingoBridge 系统日常维护与安全审计计划表}
  \label{tab:maintenance_schedule}
  \begin{tabularx}{\textwidth}{lXX}
    \toprule
    维护项目名称 & 执行时间周期 & 具体核心执行动作与责任边界 \\
    \midrule
    系统健康监视 & 每日连续实时监控 & 通过监控系统轮询健康检查接口、异常报警即时处置 \\
    数据库自动备份 & 每日凌晨自动执行 & 导出全局加密 SQL 压缩包、多副本异地冗余同步 \\
    垃圾数据清理 & 每周定期执行 & 清理过期的模拟测试数据、清理临时文件与上传缓存 \\
    隐私合规审计 & 每季度深度审计 & 核查用户自主删除记录、审计管理员及教师的操作日志 \\
    容灾恢复演练 & 每半年模拟演练 & 在独立沙盒环境中还原备份数据、检验应急灾备时效 \\
    \bottomrule
  \end{tabularx}
\end{table}

通过如表\ref{tab:maintenance_schedule}所示的细颗粒度日常维护以及严格的隐私安全保护标准，LingoBridge 平台在实现高效拼音跟读教学的同时，也筑牢了坚实的数据安全屏障，真正做到了让教师安心授课、让学生放心练习。
